\documentclass[conference]{IEEEtran}
\usepackage{amsmath,graphicx,booktabs,url}
\title{Native and Newton-Completed GridSFM Responses to Independent Active and Reactive Loading}
\author{\IEEEauthorblockN{Author list pending confirmation}}
\begin{document}
\maketitle
\begin{abstract}
This working skeleton asks which errors of frozen GridSFM persist after the released Newton power-flow completion under matched active and reactive loading. Results and contribution statements remain pending the extension audit and execution.
\end{abstract}
\begin{IEEEkeywords}Optimal power flow, power flow, neural networks, physical feasibility\end{IEEEkeywords}
\section{Introduction and Related Work}
The central question separates native prediction errors, errors in prescribed controls after completion, and operational constraint violations. Generic accuracy--feasibility differences are established knowledge; an additional contribution requires a matched, literature-bounded observation.
\section{Matched Study and Official Workflow}
Reuse the frozen three-network, 59-input protocol. Distinguish native outputs, frozen controls plus official completion, verified target-adapted controls plus completion, and downstream procedures outside this scope.
\section{Baseline and Independent P/Q Loading}
Pending validated stage integration and complete matched evaluation.
\section{Case-3 Electrical Accounting}
Pending matched native/completed distribution and equipment comparison.
\section{Focused Response Sensitivity}
Pending endpoint reuse; separate baseline offsets from direction and magnitude errors.
\section{Supported Understanding and Limitations}
No novelty, model benefit, speedup, objective-quality or generalization claim is established by this skeleton.
\bibliographystyle{IEEEtran}
\bibliography{references}
\end{document}
